\chapter{LLMs as Dynamical Systems}

\section{Introduction}

Large language models (LLMs) are, at bottom, functions: they accept a sequence of tokens and produce a probability distribution over the next token. When equipped with a decoding strategy---greedy, beam search, nucleus sampling---they become maps on the space of token sequences. When their outputs are fed back as inputs, they become discrete dynamical systems.

This chapter develops this perspective systematically. We begin by formalizing what an LLM computes as a map on symbolic sequences, connecting to the sliding block codes and subshifts of Chapters 2--3. We then study the dynamics of iterated LLM generation---feeding outputs back as inputs---and show how recent empirical findings (convergence to periodic orbits, model collapse, self-improvement) fit naturally into the framework of discrete dynamical systems developed throughout this book.

The key conceptual contribution is the \textbf{(f, x) decomposition}: the observation that a prompt simultaneously encodes a function specification $f$ and an input $x$, and that iteration can mutate both. This framework unifies phenomena as diverse as model collapse, self-refinement, and multi-agent debate under a single dynamical systems umbrella.



\section{The LLM as a Map}

\subsection{14.2.1 The Conditional Distribution}

Let $\Sigma$ be a finite token vocabulary with $|\Sigma| = V$ (typically $V \approx 32{,}000$ to $128{,}000$ for modern LLMs). An autoregressive language model defines, for each finite sequence $w = (x_1, x_2, \ldots, x_n) \in \Sigma^*$, a conditional probability distribution over the next token:

\[P_\theta(x_{n+1} \mid x_1, x_2, \ldots, x_n),\]

where $\theta$ denotes the model parameters. The full joint distribution over sequences factorizes autoregressively:

\[P_\theta(x_1, x_2, \ldots, x_N) = \prod_{n=1}^{N} P_\theta(x_n \mid x_1, \ldots, x_{n-1}).\]

This is a standard probability chain rule; the content of the model lies in the parametric form of $P_\theta$.

\subsection{14.2.2 Context Windows and Finite Memory}

In practice, transformer-based LLMs have a finite \textbf{context window} of length $k$ (typically $k = 2{,}048$ to $128{,}000$ tokens). The conditional distribution depends only on the most recent $k$ tokens:

\[P_\theta(x_{n+1} \mid x_1, \ldots, x_n) = P_\theta(x_{n+1} \mid x_{n-k+1}, \ldots, x_n)\]

for all $n \geq k$.

\begin{definition}[LLM local rule]
An LLM with context window $k$ and vocabulary $\Sigma$ defines a function

\[p: \Sigma^k \to \Delta(\Sigma),\]

where $\Delta(\Sigma)$ is the probability simplex over $\Sigma$. The function $p$ maps a $k$-tuple of tokens to a probability distribution over the next token.

This is the stochastic analogue of the block maps from Chapter 3. Specifically, if we fix a deterministic decoding strategy, we obtain a deterministic local rule.
\end{definition}

\subsection{14.2.3 Deterministic Decoding as a Sliding Block Code}

\begin{definition}[Greedy decoding map]
Given an LLM local rule $p: \Sigma^k \to \Delta(\Sigma)$, the \textbf{greedy decoding map} is the function $f: \Sigma^k \to \Sigma$ defined by

\[f(x_{n-k+1}, \ldots, x_n) = \arg\max_{s \in \Sigma} \; p(s \mid x_{n-k+1}, \ldots, x_n).\]

(Ties are broken by a fixed ordering on $\Sigma$.)
\end{definition}

\begin{proposition}
An LLM with context window $k$ and greedy decoding induces a one-sided sliding block code $\Phi: \Sigma^{\mathbb{N}} \to \Sigma^{\mathbb{N}}$ with memory $k-1$ and anticipation $0$.


\end{proposition}

\begin{proof}
The output at position $n$ depends on positions $n-k+1, \ldots, n$ of the input. This is exactly a sliding block code with memory $m = k - 1$ and anticipation $a = 0$, in the sense of Definition 3.2.2. The one-sidedness reflects the causal (autoregressive) structure: the LLM cannot look ahead. 

\begin{remark}
This is an unusual sliding block code in two respects. First, the local rule $f$ is not given by a simple formula but by a neural network with billions of parameters. Second, the code is one-sided: it acts on $\Sigma^{\mathbb{N}}$ rather than $\Sigma^{\mathbb{Z}}$, reflecting the causal nature of language generation.
\end{remark}

\end{proof}

\subsection{14.2.4 The Generated Subshift}

\begin{definition}[Generated language]
Given an LLM with greedy decoding map $f: \Sigma^k \to \Sigma$ and a set of seed prompts $\mathcal{P} \subseteq \Sigma^k$, the \textbf{generated language} is

\[\mathcal{L}(f, \mathcal{P}) = \{ w \in \Sigma^* : w \text{ is a contiguous block of some sequence generated by iterating } f \text{ from some seed in } \mathcal{P} \} .\]

The closure of all bi-infinite sequences whose every finite subword lies in $\mathcal{L}(f, \mathcal{P})$ forms a subshift $X_f \subseteq \Sigma^{\mathbb{Z}}$.
\end{definition}

\begin{example}[Tiny vocabulary]
Let $\Sigma = \{ a, b\} $, context window $k = 2$, and suppose the greedy decoding map is:

\begin{center}
\begin{tabular}{cc}
\toprule
$(x_1, x_2)$ & $f(x_1, x_2)$ \\
\midrule
$(a, a)$ & $b$ \\
$(a, b)$ & $a$ \\
$(b, a)$ & $b$ \\
$(b, b)$ & $a$ \\
\bottomrule
\end{tabular}
\end{center}

Starting from the seed $aa$: we generate $b$, giving $aab$. The window slides to $ab$, generating $a$: $aaba$. Then $ba \to b$: $aabab$. Then $ab \to a$: $aababa$. The pattern $\ldots abababab\ldots$ emerges: a period-2 orbit of the sliding window.

This is a subshift of finite type (an edge shift on the de Bruijn graph of order 2), and the attractor is the single periodic orbit $\overline{ab}$.
\end{example}

\subsection{14.2.5 Stochastic Generation as a Measure}

With temperature $\tau > 0$, the LLM generates tokens by sampling from the softmax distribution:

\[P_\tau(x_{n+1} = s \mid x_{n-k+1}, \ldots, x_n) = \frac{\exp(\ell_s / \tau)}{\sum_{s' \in \Sigma} \exp(\ell_{s'} / \tau)},\]

where $\ell_s$ are the logits. As $\tau \to 0$, this concentrates on the greedy argmax. As $\tau \to \infty$, it approaches the uniform distribution.

The stochastic process defined by sampling with a fixed temperature defines a \textbf{shift-invariant probability measure} $\mu_\tau$ on $\Sigma^{\mathbb{N}}$, supported on a subshift that depends on $\tau$. In the language of ergodic theory, $(\Sigma^{\mathbb{N}}, \sigma, \mu_\tau)$ is a measure-preserving system (after suitable stationarization), and the entropy rate of the process is

\[h(\mu_\tau) = -\lim_{n \to \infty} \frac{1}{n} \sum_{w \in \Sigma^n} \mu_\tau([w]) \log \mu_\tau([w]),\]

where $[w]$ is the cylinder set of sequences beginning with $w$. This is the Shannon entropy rate of the LLM's output process.



\section{Iterated LLM Generation}

\subsection{14.3.1 The Iteration Scheme}

We now study the core dynamical system: feeding an LLM's output back as its input.

\begin{definition}[LLM iteration]
Let $\mathrm{LLM}: \Sigma^\textit{ \to \Sigma^}$ denote the map that takes a prompt string to the LLM's full response (a finite string produced by generating tokens until a stop condition). Define the \textbf{iteration}:

\[x_0 \in \Sigma^*, \qquad x_{n+1} = \mathrm{LLM}(x_n), \qquad n = 0, 1, 2, \ldots\]

This is a discrete dynamical system $(S, F)$ where $S = \Sigma^*$ (finite strings over the token vocabulary) and $F = \mathrm{LLM}$.
\end{definition}

\begin{remark}
There is a subtlety: $\Sigma^*$ is countably infinite, not finite. The theory of Chapter 1 (iterated maps on finite sets) does not directly apply---orbits need not be eventually periodic. However, if we restrict attention to strings of bounded length (as enforced by a maximum output length), the state space becomes finite and all orbits are eventually periodic.
\end{remark}

\subsection{14.3.2 Convergence to Periodic Attractors}

\begin{theorem}[Wang et al., 2025]
In experiments with successive paraphrasing---where an LLM is asked to rephrase its own output---the dynamics converge to \textbf{2-cycles} (period-2 attractors). That is, there exist strings $a, b \in \Sigma^*$ with $\mathrm{LLM}(a) = b$ and $\mathrm{LLM}(b) = a$, and orbits from generic initial conditions converge to one of these 2-cycles.

\textit{Discussion.} This is Class 2 behavior in the Wolfram classification (Chapter 7): the system converges to periodic attractors from generic initial conditions. The 2-cycles arise because the LLM has systematic biases---preferred phrasings, syntactic patterns, and lexical choices---that act as attractors under iteration. A paraphrase of a paraphrase tends to oscillate between two "canonical forms" that are paraphrases of each other.
\end{theorem}

\begin{example}
Consider successive paraphrasing of "The cat sat on the mat." A typical trajectory might look like:

- $x_0$: "The cat sat on the mat."
- $x_1$: "A feline was seated upon the mat."
- $x_2$: "A cat was sitting on the mat."
- $x_3$: "A feline was seated on the mat."
- $x_4$: "A cat was sitting on the mat."

From $x_2$ onward, the system oscillates between two phrasings: a 2-cycle.
\end{example}

\subsection{14.3.3 Geometric Framework}

Tacheny (2025) provides a formal geometric framework for LLM iteration, identifying two regimes:

\begin{enumerate}
\item \textbf{Contractive regime.} The iteration map is (locally) contractive in an appropriate metric on string space. Orbits converge to fixed points or short cycles. This occurs with "neutral" rewriting tasks like paraphrasing or summarization.
\end{enumerate}

\begin{enumerate}
\item \textbf{Exploratory regime.} The iteration map is expansive: small perturbations grow. Orbits diverge and exhibit sensitive dependence on initial conditions. This occurs with "adversarial" tasks like "summarize the following text, then negate the main claim."
\end{enumerate}

\begin{proposition}[Informal]
The choice of prompt---specifically, the instruction component---controls which regime the iteration inhabits. Prompts that preserve meaning (paraphrase, reformat) tend to produce contractive dynamics. Prompts that transform meaning (negate, extrapolate) tend to produce expansive dynamics.

This is a direct analogue of the role of the local rule in cellular automata (Chapter 7): different rules produce different Wolfram classes.
\end{proposition}

\subsection{14.3.4 The Telephone Game}

Perez et al. (2024) study iterated LLM generation as a "cultural telephone game": a message is passed from LLM to LLM (or from the same LLM to itself), and small biases accumulate across iterations. The key finding is that outputs converge toward \textbf{cultural attractors}---stereotypical, high-frequency patterns in the training data.

This is the dynamical systems analogue of the \textbf{basin of attraction} (Chapter 1, Definition 1.8): the set of initial conditions that converge to a given attractor. The cultural attractors of an LLM are the periodic orbits or fixed points of the iteration map, and their basins of attraction encompass the space of "reasonable" inputs.



\section{The (f, x) Decomposition}

\subsection{14.4.1 Prompts Encode Functions and Inputs}

The key structural observation underlying all of iterated LLM generation is:

\begin{observation}
A prompt $w \in \Sigma^*$ simultaneously encodes two components:
- A \textbf{function specification} $f$: what the LLM should do (the instruction).
- An \textbf{input} $x$: what the LLM should do it to (the data).

The LLM's output is (approximately) $f(x)$.
\end{observation}

\begin{example}
Consider the prompt:

> "Given x=3, add one and output in the same format: Given x=[result], add one and output in the same format."

This encodes:
- $f$: the function "add one to $x$ and format the result as a new prompt with the same instruction."
- $x = 3$: the current input.

The expected output is:

> "Given x=4, add one and output in the same format: Given x=[result], add one and output in the same format."

This output encodes the same function $f$ but with updated input $x' = 4 = f(3)$.
\end{example}

\subsection{14.4.2 The (f, x) Dynamical System}

\begin{definition}[Function-input decomposition]
Let $\mathcal{F}$ denote the space of "function specifications" (the instruction component of a prompt) and $\mathcal{X}$ the space of "inputs" (the data component). An LLM iteration defines a map

\[\Psi: \mathcal{F} \times \mathcal{X} \to \mathcal{F} \times \mathcal{X}\]

given by

\[\Psi(f, x) = (\varphi(f, x), \; f(x)),\]

where:
\begin{itemize}
\item $f(x)$ is the result of applying the function $f$ to the input $x$ (the "computation" component).
\item $\varphi(f, x)$ is the \textbf{meta-rule}: the function specification that appears in the output (the "instruction mutation" component).

\end{itemize}
\textbf{Case 1: Perfect instruction preservation ($\varphi(f, x) = f$ for all $x$).} The dynamics reduce to ordinary iteration of $f$:

\[(f, x) \to (f, f(x)) \to (f, f^2(x)) \to \cdots\]

This is the setting of Chapter 1: a fixed map iterated on a state space.

\textbf{Case 2: Instruction mutation ($\varphi \neq \mathrm{id}$).} Both the function and the input evolve:

\[(f_0, x_0) \to (f_1, x_1) \to (f_2, x_2) \to \cdots\]

where $f_{n+1} = \varphi(f_n, x_n)$ and $x_{n+1} = f_n(x_n)$. This is a \textbf{self-modifying dynamical system}: the rule itself changes at each step.
\end{definition}

\begin{remark}
Self-modifying dynamics are well-studied in the context of \textbf{iterated function systems} (Chapter 6), where the map applied at each step is chosen from a finite set. The (f, x) framework generalizes this by allowing the map to mutate continuously.
\end{remark}

\subsection{14.4.3 Worked Example: Iterated Increment}

Let $\mathcal{X} = \mathbb{Z}$, $f_0(x) = x + 1$, and suppose the LLM preserves instructions perfectly: $\varphi(f, x) = f$.

Then:
\begin{itemize}
\item $(f_0, 3) \to (f_0, 4) \to (f_0, 5) \to \cdots$
\end{itemize}

The orbit is $\{ 3, 4, 5, \ldots\} $---a divergent trajectory with no periodic orbit (since $\mathcal{X}$ is infinite). If $\mathcal{X}$ is taken modulo $N$ (i.e., $\mathcal{X} = \mathbb{Z}/N\mathbb{Z}$), the orbit is periodic with period $N$.

\subsection{14.4.4 Worked Example: Instruction Drift}

Now suppose the LLM is imperfect: it occasionally mutates the instruction. Let $f_0(x) = x + 1$ but with probability $\epsilon$ per iteration the instruction drifts to $f'(x) = x + 2$, and once there, it stays.

The trajectory might look like:
\begin{itemize}
\item $(f_0, 3) \to (f_0, 4) \to (f', 5) \to (f', 7) \to (f', 9) \to \cdots$
\end{itemize}

The "increment by 1" rule mutated to "increment by 2" at step 2, and all subsequent dynamics follow $f'$. This illustrates how instruction mutation changes the qualitative behavior of the system.



\section{Model Collapse}

\subsection{14.5.1 The Phenomenon}

Model collapse occurs when a generative model is trained on data produced by itself (or by a previous generation of itself). The distribution of outputs progressively loses diversity and converges to a low-variance, low-entropy approximation of the original.

\begin{theorem}[Shumailov et al., 2024]
Consider a sequence of models $M_0, M_1, M_2, \ldots$ where $M_0$ is trained on real data, and $M_{n+1}$ is trained on data generated by $M_n$. Then:
1. \textbf{Tail erosion.} The tails of the learned distribution shrink at each generation. Rare events become rarer and eventually disappear.
2. \textbf{Variance collapse.} For Gaussian models with training set size $N$, the variance at generation $n$ satisfies

\[\sigma_n^2 \approx \sigma_0^2 \cdot \frac{1}{1 + n/N}.\]

The variance decreases as $O(1/n)$.

3. \textbf{Convergence to a point mass.} In the limit $n \to \infty$, the model's distribution collapses to a point mass (or near-point-mass) at the mode.
\end{theorem}

\subsection{14.5.2 Information-Theoretic Explanation}

\begin{proposition}[Data Processing Inequality]
Let $X$ be the true data distribution, $M_n$ the model at generation $n$, and $Y_n$ the data generated by $M_n$. By the data processing inequality,

\[I(X; Y_{n+1}) \leq I(X; Y_n),\]

where $I(\cdot; \cdot)$ denotes mutual information. Each generation of self-training is a lossy channel: information about the original distribution can only decrease.


\end{proposition}

\begin{proof}
The chain $X \to Y_n \to M_{n+1} \to Y_{n+1}$ is a Markov chain. The data processing inequality states that $I(X; Y_{n+1}) \leq I(X; Y_n)$ for any Markov chain $X \to Y \to Z$. 

\begin{corollary}
Self-training cannot, in expectation, increase the fidelity of a model to the true data distribution. Model collapse is thermodynamically inevitable in the absence of fresh data.
\end{corollary}

\end{proof}

\subsection{14.5.3 Self-Consuming Loops}

Alemohammad et al. (2023) identify three types of \textbf{self-consuming generative loops}:

\begin{enumerate}
\item \textbf{Fully synthetic loop.} All training data for $M_{n+1}$ comes from $M_n$. Collapse is fastest.
\item \textbf{Synthetic augmentation loop.} Training data is a mixture of real and synthetic data. Collapse is slower but still occurs unless the fraction of real data exceeds a threshold.
\item \textbf{Fresh data loop.} Each generation receives fresh real data. Collapse is avoided.
\end{enumerate}

They term the pathological behavior \textbf{Model Autophagy Disorder (MAD)}: the model consumes itself and degenerates.

\subsection{14.5.4 Model Collapse in the (f, x) Framework}

In the (f, x) decomposition, model collapse corresponds to a \textbf{destructive meta-rule} $\varphi$:

\[\varphi(f, x) = f_{\text{degraded}},\]

where $f_{\text{degraded}}$ is a lower-fidelity version of $f$. At each generation:

\[(f_0, x_0) \to (f_1, f_0(x_0)) \to (f_2, f_1(f_0(x_0))) \to \cdots\]

with $f_0 \succ f_1 \succ f_2 \succ \cdots$ in some quality ordering. The sequence converges to $(f_\infty, x_\infty)$ where $f_\infty$ is maximally degenerate (e.g., it maps all inputs to the mode) and $x_\infty$ is a fixed point of $f_\infty$.

\begin{example}
Consider a Gaussian model. Let $f_n$ denote "sample from $\mathcal{N}(\mu, \sigma_n^2)$." Each self-training step shrinks $\sigma_n$:

\[\sigma_{n+1}^2 = \sigma_n^2 \cdot \frac{N}{N + 1} < \sigma_n^2.\]

The fixed point is $\sigma_\infty^2 = 0$: the collapsed model is a delta function at $\mu$.
\end{example}



\section{Self-Improvement and Its Limits}

\subsection{14.6.1 Self-Refine}

Madaan et al. (2023) propose \textbf{Self-Refine}: a three-step iterative loop.

\begin{enumerate}
\item \textbf{Generate.} The LLM produces an initial response $y_0 = \mathrm{LLM}(x)$.
\item \textbf{Critique.} The LLM evaluates its own response: $c_0 = \mathrm{LLM}(\text{"Critique: "} \| y_0)$.
\item \textbf{Refine.} The LLM improves the response using the critique: $y_1 = \mathrm{LLM}(\text{"Refine: "} \| y_0 \| c_0)$.
\end{enumerate}

Steps 2--3 repeat: $y_0 \to y_1 \to y_2 \to \cdots$.

\begin{observation}
Empirically, the quality of $y_n$ improves for the first few iterations and then \textbf{saturates}: $y_n \approx y_{n+1}$ for $n$ sufficiently large. This is convergence to an \textbf{approximate fixed point} of the refine map.

In the (f, x) framework: $f$ = "refine the response" and $x_n = y_n$. The dynamics are:

\[x_{n+1} = f(x_n),\]

and the system converges to a fixed point $x^\textit{$ with $f(x^}) \approx x^*$. The convergence is consistent with the refine map being \textbf{contractive} in some implicit quality metric.
\end{observation}

\subsection{14.6.2 Self-Rewarding}

Yuan et al. (2024) propose a more ambitious scheme: the LLM serves simultaneously as \textbf{policy} (generator) and \textbf{reward model} (evaluator). At each iteration:

\begin{enumerate}
\item The LLM generates responses.
\item The same LLM scores them (acting as reward model).
\item The LLM is fine-tuned on high-scoring responses (reinforcement learning from self-feedback).
\end{enumerate}

This is a genuine (f, x) system where \textbf{both components improve}:

\[\Psi(f_n, x_n) = (f_{n+1}, x_{n+1})\]

with $f_{n+1}$ being a better policy/reward model than $f_n$, and $x_{n+1}$ being a higher-quality response than $x_n$. Empirically, self-rewarding models improve over multiple iterations, with each generation outperforming the last on benchmarks.

The contrast with model collapse (Section 14.5) is instructive: here $\varphi$ is \textit{constructive} (improving $f$) rather than destructive (degrading $f$). The difference lies in the feedback mechanism: self-rewarding includes an explicit quality signal, while self-consuming loops do not.

\subsection{14.6.3 Limits of Self-Correction}

\begin{theorem}[Huang et al., 2023, informal]
LLMs cannot reliably self-correct their reasoning without external feedback. Specifically: for tasks requiring multi-step logical reasoning, the "critique-and-revise" map often \textbf{worsens} correct answers---the refinement map has wrong fixed points.

\textit{Discussion.} The issue is that the LLM's internal model of "correctness" is imperfect. When the same flawed model is used for both generation and evaluation, systematic errors reinforce themselves. In dynamical systems terms: the refinement map has \textbf{spurious fixed points} (wrong answers that the model believes are correct), and the basin of attraction of these spurious fixed points can be large.

This is analogous to the fixed-point structure of Chapter 1: a map can have both "good" fixed points (correct answers) and "bad" fixed points (wrong answers that are stable under iteration). Without external perturbation, the system cannot escape a bad basin.
\end{theorem}

\subsection{14.6.4 Constitutional AI}

Bai et al. (2022) introduce \textbf{Constitutional AI (CAI)}: a set of explicit principles ("the constitution") that constrain the LLM's behavior during self-improvement. In the (f, x) framework, the constitution acts as a \textbf{constraint on the meta-rule} $\varphi$:

\[\varphi(f, x) \in \mathcal{C},\]

where $\mathcal{C} \subset \mathcal{F}$ is the set of function specifications consistent with the constitution. This prevents $\varphi$ from drifting toward degenerate or harmful behaviors---it restricts the dynamics to a controlled subspace.



\section{Computational Power of Iterated LLMs}

\subsection{14.7.1 A Single Forward Pass: TC$^0$}

\begin{theorem}[Merrill and Sabharwal, 2024]
A single forward pass of a transformer with $O(1)$ layers, $O(\mathrm{poly}(n))$ embedding dimension, and $O(1)$ precision per parameter lies in the circuit complexity class $\mathrm{TC}^0$: the class of problems solvable by constant-depth, polynomial-size threshold circuits.

$\mathrm{TC}^0$ is powerful---it contains addition, multiplication, sorting, and matrix operations---but it does not contain all of $\mathrm{P}$. In particular, problems requiring $\omega(1)$-depth circuits (such as evaluating Boolean formulas of unbounded depth) cannot be solved in a single forward pass.
\end{theorem}

\subsection{14.7.2 Chain of Thought: Turing Completeness}

\begin{theorem}[Merrill and Sabharwal, 2024]
A transformer augmented with $O(\mathrm{poly}(n))$ chain-of-thought (CoT) steps---where intermediate tokens are generated and fed back---can simulate any polynomial-time Turing machine. That is, transformers with CoT are \textbf{Turing complete} (up to the polynomial bound on the number of steps).

\textit{Proof sketch.} Each CoT step is a $\mathrm{TC}^0$ computation. A polynomial number of such steps can simulate a polynomial-time Turing machine step by step, using the CoT tokens as the tape. The transformer encodes the transition function, and each forward pass executes one (or a bounded number of) steps. 

\end{theorem}


\subsection{14.7.3 Prompt-Based Turing Completeness}

\begin{theorem}[Qiu et al., 2024]
For any computable function $g: \{ 0,1\} ^\textit{ \to \{ 0,1\} ^}$, there exists a prompt $w$ such that a fixed transformer, given $w$ concatenated with input $x$, computes $g(x)$ using chain-of-thought. That is, the space of prompts parametrizes all computable functions.

This is a remarkable result: \textbf{prompting alone suffices for Turing completeness.} The prompt plays the role of a program, and the transformer plays the role of a universal machine.
\end{theorem}

\subsection{14.7.4 Implications for the (f, x) Framework}

Combining Theorems 14.20 and 14.21 with the (f, x) decomposition:

\begin{corollary}
The dynamical system $\Psi: (f, x) \mapsto (\varphi(f, x), f(x))$, instantiated with a sufficiently powerful LLM and chain-of-thought, can in principle realize any computable dynamical system. In particular, for any computable map $g: \mathcal{X} \to \mathcal{X}$, there exists a prompt encoding $f = g$ such that the iteration $x_{n+1} = f(x_n)$ computes the orbit of $g$.

This means the theory of LLM iteration inherits the full richness (and undecidability) of general computation. Questions like "does this iterated LLM converge?" are, in the general case, undecidable (by reduction to the halting problem).
\end{corollary}



\section{Multi-Agent LLM Systems}

\subsection{14.8.1 The Product Dynamical System}

When multiple LLM agents interact, the state space becomes a product. Let agents $A$ and $B$ maintain internal states (conversation histories) $s_A, s_B \in \Sigma^\textit{$, and let $m \in \Sigma^}$ be the current message being exchanged.

\begin{definition}[Two-agent LLM system]
A \textbf{coupled LLM dynamical system} with agents $A, B$ is a map

\[\Psi: \Sigma^\textit{ \times \Sigma^} \times \Sigma^\textit{ \to \Sigma^} \times \Sigma^\textit{ \times \Sigma^}\]

defined by alternating updates:

\[\Psi(s_A, s_B, m) = \begin{cases} (s_A \| m, \; s_B, \; \mathrm{LLM}_A(s_A \| m)) \& \text{if it is } A\text{'s turn,} \\ (s_A, \; s_B \| m, \; \mathrm{LLM}_B(s_B \| m)) \& \text{if it is } B\text{'s turn,} \end{cases}\]

where $\|$ denotes concatenation.
\end{definition}

\subsection{14.8.2 Role-Playing Agents: CAMEL}

Li et al. (2023) instantiate this framework with \textbf{CAMEL}: two LLM agents assigned complementary roles (e.g., "AI user" and "AI assistant") that collaborate through structured dialogue. Each agent's prompt includes a role specification and a task description.

In the (f, x) framework, this is:

\[\Psi(f_A, f_B, x) = (f_A, f_B, f_B(f_A(x))),\]

where $f_A$ is the "user" role (generates requests) and $f_B$ is the "assistant" role (generates responses). The two roles are fixed, and only $x$ (the conversation content) evolves. This is a composition: the effective single-step map is $f_B \circ f_A$.

\subsection{14.8.3 Emergent Collective Dynamics}

Park et al. (2023) simulate \textbf{25 generative agents} in a virtual town, each with:
\begin{itemize}
\item A memory stream (list of past observations and reflections).
\item A retrieval mechanism (attention over memory).
\item A planning module (daily schedule generation).
\item A reflection module (periodic synthesis of memories into higher-level insights).
\end{itemize}

The resulting multi-agent system exhibits emergent social behaviors: agents form relationships, spread information, coordinate activities, and throw a party that was not explicitly programmed.

In dynamical systems terms, the state is a vector $(s_1, s_2, \ldots, s_{25}) \in (\Sigma^*)^{25}$, and the update map couples all agents through a shared environment. The emergent social structures---stable relationships, group activities, information cascades---are \textbf{collective attractors} of this high-dimensional coupled dynamical system. They are analogous to the synchronized oscillations and pattern formation seen in coupled map lattices (a continuous-space analogue of the coupled cellular automata of Chapter 7).



\section{Open Questions}

The dynamical systems perspective on LLMs raises several fundamental questions, many of which connect back to themes developed throughout this book.

\begin{question}[Attractor control]
Can we design meta-rules $\varphi$ that avoid both model collapse (convergence to degenerate fixed points) and chaos (no convergence)? What is the analogue of "edge of chaos" for LLM iteration---the boundary between contractive dynamics (Section 14.3.3, regime 1) and expansive dynamics (regime 2)?

In the cellular automata setting (Chapter 7), the edge of chaos corresponds to Wolfram Class 4 rules, which produce complex, long-lived transients and support computation. The analogous question for LLM iteration is whether there exist prompts that induce Class 4 behavior: complex, structured, non-repeating output that nonetheless has internal organization.
\end{question}

\begin{question}[Logical depth of LLM outputs]
Is there a formal connection between the \textbf{logical depth} (Chapter 9) of an LLM's output and the "creativity" or "interestingness" of the generation process? Recall that logical depth measures the computational effort required to produce a string from its shortest description. An LLM that generates logically deep text is, in some sense, performing substantial computation---not producing trivial repetition (shallow) or random noise (incompressible but also shallow in Bennett's framework).
\end{question}

\begin{question}[Prompt engineering as dynamical systems design]
Prompt engineering is, in the (f, x) framework, the design of the initial function specification $f_0$ and the meta-rule $\varphi$. Can the attractor structure of the resulting dynamical system be \textit{predicted} from properties of the prompt? Is there a "bifurcation theory" for prompt parameters---small changes in the prompt that cause qualitative changes in the long-term dynamics?
\end{question}

\begin{question}[Fresh data as external forcing]
Alemohammad et al. (2023) show that fresh data prevents model collapse. In dynamical systems terms, fresh data is an \textbf{external forcing} or \textbf{perturbation} that prevents the system from settling into a degenerate attractor. What is the minimum rate of fresh data injection needed to maintain a given level of distributional diversity? This is analogous to the problem of \textbf{noise-sustained patterns} in nonlinear dynamics.
\end{question}

\begin{question}[Multi-agent phase transitions]
As the number of LLM agents and the topology of their interactions change, do the collective dynamics exhibit phase transitions? Is there a critical connectivity at which coherent collective behavior emerges, analogous to percolation thresholds or the onset of synchronization in coupled oscillator networks?
\end{question}



\section{Summary}

\begin{center}
\begin{tabular}{cc}
\toprule
Concept & Dynamical Systems Formalization \\
\midrule
LLM with greedy decoding & One-sided sliding block code, memory $k-1$ \\
LLM with temperature $\tau > 0$ & Shift-invariant measure $\mu_\tau$ on $\Sigma^{\mathbb{N}}$ \\
Iterated generation & DDS $(\Sigma^*, \mathrm{LLM})$ \\
Successive paraphrasing & Convergence to 2-cycles (Class 2) \\
Prompt = (instruction, data) & (f, x) decomposition \\
Perfect instruction following & Ordinary iteration: $x_{n+1} = f(x_n)$ \\
Instruction drift & Self-modifying system: $(f_{n+1}, x_{n+1}) = (\varphi(f_n, x_n), f_n(x_n))$ \\
Model collapse & Destructive $\varphi$; convergence to degenerate fixed point \\
Self-improvement & Constructive $\varphi$; convergence to improved fixed point \\
Constitutional AI & Constraint $\varphi(f, x) \in \mathcal{C}$ \\
Single forward pass & $\mathrm{TC}^0$ computation \\
With chain of thought & Turing-complete computation \\
Multi-agent system & Coupled DDS on product state space \\
\bottomrule
\end{tabular}
\end{center}

The perspective developed in this chapter reveals that many apparently disparate phenomena in LLM research---model collapse, self-improvement, multi-agent emergence, prompt sensitivity---are manifestations of a single underlying structure: the dynamics of iterated maps on the space of strings. The tools of Chapters 1--9---functional graphs, symbolic dynamics, sliding block codes, entropy, logical depth---apply directly and provide both conceptual clarity and quantitative predictions.



\section{References}

\begin{itemize}
\item Alemohammad, S., Casco-Rodriguez, J., Luzi, L., Humayun, A. I., Babaei, H., LeJeune, D., Siahkoohi, A., and Baraniuk, R. G. (2023). Self-consuming generative models go MAD. \textit{arXiv preprint} arXiv:2307.01850.
\end{itemize}

\begin{itemize}
\item Bai, Y., Kadavath, S., Kundu, S., Askell, A., Kernion, J., Jones, A., Chen, A., Goldie, A., Mirhoseini, A., McKinnon, C., et al. (2022). Constitutional AI: Harmlessness from AI feedback. \textit{arXiv preprint} arXiv:2212.08073.
\end{itemize}

\begin{itemize}
\item Huang, J., Chen, X., Mishra, S., Zheng, H. S., Yu, A. W., Song, X., and Zhou, D. (2023). Large language models cannot self-correct reasoning yet. \textit{arXiv preprint} arXiv:2310.01798.
\end{itemize}

\begin{itemize}
\item Li, G., Hammoud, H. A. A. K., Itani, H., Khizbullin, D., and Ghanem, B. (2023). CAMEL: Communicative agents for "mind" exploration of large language model society. \textit{arXiv preprint} arXiv:2303.17760.
\end{itemize}

\begin{itemize}
\item Lind, D. and Marcus, B. (2021). \textit{An Introduction to Symbolic Dynamics and Coding}. 2nd edition. Cambridge University Press.
\end{itemize}

\begin{itemize}
\item Madaan, A., Tandon, N., Gupta, P., Hallinan, S., Gao, L., Wiegreffe, S., Alon, U., Dziri, N., Prabhumoye, S., Yang, Y., et al. (2023). Self-Refine: Iterative refinement with self-feedback. \textit{arXiv preprint} arXiv:2303.17651.
\end{itemize}

\begin{itemize}
\item Merrill, W. and Sabharwal, A. (2024). The expressive power of transformers with chain of thought. \textit{arXiv preprint} arXiv:2310.07923.
\end{itemize}

\begin{itemize}
\item Park, J. S., O'Brien, J. C., Cai, C. J., Morris, M. R., Liang, P., and Bernstein, M. S. (2023). Generative agents: Interactive simulacra of human behavior. \textit{arXiv preprint} arXiv:2304.03442.
\end{itemize}

\begin{itemize}
\item Perez, C. R., Morvant, E., and Gauthier, E. (2024). Cultural attractors in LLM telephone games. \textit{Proceedings of the Annual Meeting of the Cognitive Science Society}.
\end{itemize}

\begin{itemize}
\item Qiu, L., Jiang, N., and Ma, T. (2024). Prompting is Turing-complete. \textit{arXiv preprint} arXiv:2411.01992.
\end{itemize}

\begin{itemize}
\item Shumailov, I., Shumaylov, Z., Zhao, Y., Gal, Y., Papernot, N., and Anderson, R. (2024). The curse of recursion: Training on generated data makes models forget. \textit{arXiv preprint} arXiv:2305.17493.
\end{itemize}

\begin{itemize}
\item Tacheny, N. (2025). A geometric framework for iterative LLM generation. \textit{arXiv preprint} arXiv:2512.10350.
\end{itemize}

\begin{itemize}
\item Wang, S., Dong, Y., and Cotterell, R. (2025). Successive paraphrasing converges to fixed cycles. \textit{arXiv preprint} arXiv:2502.15208.
\end{itemize}

\begin{itemize}
\item Yuan, W., Pang, R. Y., Cho, K., Sukhbaatar, S., Xu, J., and Weston, J. (2024). Self-rewarding language models. \textit{arXiv preprint} arXiv:2401.10020.
\end{itemize}



\section{Recommended Reading}

For the dynamical systems perspective on LLMs:

\begin{itemize}
\item \textbf{Wang et al. (2025)} is essential reading: the first rigorous demonstration of convergence to periodic orbits in LLM iteration. Short and well-written.
\end{itemize}

\begin{itemize}
\item \textbf{Shumailov et al. (2024)} is the definitive treatment of model collapse. The theoretical analysis is clean and the experimental evidence compelling.
\end{itemize}

For computational expressiveness of transformers:

\begin{itemize}
\item \textbf{Merrill \& Sabharwal (2024)} establishes the TC$^0$ bound for single forward passes and Turing completeness with chain-of-thought. Accessible and important.
\end{itemize}

\begin{itemize}
\item \textbf{Qiu et al. (2024)} proves Turing completeness via prompting alone. Short paper with a beautiful result.
\end{itemize}

For multi-agent LLM systems:

\begin{itemize}
\item \textbf{Park et al. (2023)} on Generative Agents is a tour de force: 25 agents exhibiting emergent social behavior. The supplementary materials detail the architecture.
\end{itemize}

\begin{itemize}
\item \textbf{Li et al. (2023)} on CAMEL provides a simpler two-agent framework useful for understanding coupled dynamics.
\end{itemize}

For self-improvement and its limits:

\begin{itemize}
\item \textbf{Madaan et al. (2023)} on Self-Refine shows what works. \textbf{Huang et al. (2023)} shows what doesn't. Reading both together gives a balanced picture.
\end{itemize}

For background on symbolic dynamics and sliding block codes:

\begin{itemize}
\item \textbf{Lind \& Marcus (2021)} is the standard reference. Chapters 1--3 and 6 are directly relevant to understanding LLMs as block maps.
\end{itemize}

For thinking about AI systems dynamically:

\begin{itemize}
\item This textbook's earlier chapters, especially Chapters 1 (iterated maps), 7 (cellular automata), and 9 (logical depth), provide the conceptual foundations for analyzing LLM iteration.
\end{itemize}
